% Generated draft: requires mathematical and source review.
\documentclass[11pt]{article}
\usepackage[margin=25mm]{geometry}
\usepackage{fontspec}
\setmainfont{Latin Modern Roman}
\usepackage{amsmath,amssymb,mathtools,unicode-math}
\setmathfont{Latin Modern Math}
\usepackage{enumitem,needspace,xurl,hyperref}
\setlength{\parindent}{0pt}
\setlength{\parskip}{5pt}
\setlength{\emergencystretch}{3em}
\newcommand{\N}{\mathbb{N}}
\newcommand{\Z}{\mathbb{Z}}
\newcommand{\Q}{\mathbb{Q}}
\newcommand{\R}{\mathbb{R}}
\newcommand{\C}{\mathbb{C}}
\newcommand{\F}{\mathbb{F}}
\newcommand{\A}{\mathbb{A}}
\newcommand{\PP}{\mathbb{P}}
\newcommand{\E}{\mathbb{E}}
\newcommand{\eps}{\varepsilon}
\newcommand{\dd}{\,\mathrm{d}}
\begin{document}
\section*{UnsolvedMath ID: 801108}
\textbf{Source problem code:} catalogue entry
\subsection{Definitions and mathematical statement}
For a discrete group \(H\), let \(L(H)=\{\lambda_H(h):h\in H\}''\subseteq B(\ell^2(H))\) be its group von Neumann algebra, where \((\lambda_H(h)\xi)(k)=\xi(h^{-1}k)\). Let \(\mathbb F_n\) denote the free group of rank \(n\). Here \(\cong\) denotes isomorphism of von Neumann algebras.

Under the usual model-theoretic interpretation, an existential inclusion of \(\mathbb F_2\) in \(G\) is an injective group homomorphism \(i:\mathbb F_2\hookrightarrow G\) such that, for every quantifier-free formula \(\varphi(\bar x,\bar a)\) in the language of groups, with finite parameter tuple \(\bar a\) from \(\mathbb F_2\),
\[
\mathbb F_2\models\exists\bar x\,\varphi(\bar x,\bar a)
\quad\Longleftrightarrow\quad
G\models\exists\bar x\,\varphi(\bar x,i(\bar a)).
\]
This interpretation of ``existential inclusion'' is not explicitly defined in the supplied record.

The literal first question is whether
\[
\exists\,G\;\exists\,i:\mathbb F_2\hookrightarrow G\;\exists\,n\in\mathbb N
\quad\text{such that}\quad
i\text{ is an existential inclusion and }L(G)\not\cong L(\mathbb F_n).
\]
The convention for whether \(0\in\mathbb N\) is not supplied. No countability or factor hypothesis on \(G\) is explicitly stated.

The source then asks, specifically, whether there exists a ``non elementary surface group'' \(G\) such that \(L(G)\not\cong L(\mathbb F_n)\). The surface class and the meaning of ``non elementary'' are not defined in the supplied record. It is also not explicit whether this second question retains the existential-inclusion requirement or how its rank \(n\) is quantified.

\emph{Unresolved qualification.} The phrase ``for some \(n\in\mathbb N\)'' in the first question gives the existential rank quantifier displayed above. Requiring instead that \(L(G)\not\cong L(\mathbb F_n)\) for every finite rank \(n\) would be a different formulation. The supplied record does not justify replacing the literal existential quantifier by that universal quantifier. These issues require source clarification before a single precise intended formulation can be fixed.
\subsection{Short English statement}
The literal question asks for a group containing an existentially embedded free group of rank two whose group von Neumann algebra differs from that of a free group of some finite rank. A further question asks for a non elementary surface group with such a nonisomorphism. The intended rank quantifier, the surface-group class, and the hypotheses inherited by the second question require clarification.
\subsection{Sources}
\begin{itemize}
\item Supplied local UnsolvedMath registry record: permanent ID \(801108\), source problem code \texttt{catalogue entry}, title ``Ultraproduct \(\mathrm{II}_1\) Factors --- Problem 9.'' Upstream revision: \texttt{d371ee917ac9647fbb06cd6fc529deec6f3dda82}.
\item Attribution supplied by the registry: Elayavalli, proposed in 2025; source list ``Elayavalli --- 50 Open Problems: Ultraproduct \(\mathrm{II}_1\) Factors (2025),'' Problem 9. Supplied source URL: \url{https://arxiv.org/abs/2511.20377}.
\item The supplied provenance records access on 2026-10-16 and extraction by \texttt{arxiv\_latex}. These are supplied provenance claims; the source has not been independently accessed or verified for this draft.
\item The registry labels this problem ``open.'' This is a supplied status label, not an independently verified assessment of the literature.
\end{itemize}
\end{document}
